\begin{IEEEbiography}[{\includegraphics[width=1in,height=1.25in,clip,keepaspectratio]{images/author_ricky.jpg}}]{Ricky Eka Putra}
menyelesaikan studi Sarjana, Magister, dan Doktor di bidang Ilmu Komputer di Institut Teknologi Sepuluh Nopember (ITS) masing-masing pada tahun 2008, 2011, dan 2020, serta memperoleh gelar Profesi Insinyur pada tahun 2024. Saat ini, beliau bertugas sebagai dosen dan peneliti di Departemen Informatika, Fakultas Informatika, Universitas Negeri Surabaya. Beliau aktif berkiprah dalam pelaksanaan tridarma perguruan tinggi yang mencakup pengajaran bidang informatika, program pemberdayaan masyarakat, publikasi karya ilmiah pada berbagai jurnal bereputasi, diseminasi riset pada forum seminar ilmiah, serta perolehan Hak Kekayaan Intelektual (HKI). Fokus penelitian utama beliau diarahkan pada pengembangan model kecerdasan buatan terapan yang menjembatani kemajuan metode komputasi cerdas dengan penyelesaian tantangan berbasis data nyata. Bidang minat penelitian beliau meliputi Deep Learning, Computer Vision, Big Data Analytics, aplikasi digital dan user experience, psikologi klinis, intensi perilaku, healthcare informatics, metode statistika, serta pendidikan digital dan anak usia dini. Profil akademik dan rekaman publikasi beliau dapat diakses melalui tautan \href{https://orcid.org/0000-0002-5515-7967}{ORCID: 0000-0002-5515-7967}.
\end{IEEEbiography}

\begin{IEEEbiography}[{\includegraphics[width=1in,height=1.25in,clip,keepaspectratio]{images/author_rezky.jpg}}]{Rezky Arisanti Putri}
meraih gelar Sarjana Teknik Informatika pada tahun 2019 dengan peminatan bidang Kecerdasan Buatan (Artificial Intelligence) dan telah menuntaskan studi Magister Informatika di Universitas Negeri Surabaya (UNESA). Saat ini, beliau aktif berkiprah sebagai peneliti di Departemen Informatika, Fakultas Informatika, Universitas Negeri Surabaya. Beliau terlibat intensif dalam riset pengenalan atribut demografis wajah berbasis model transformer serta secara aktif mendiseminasikan kontribusi keilmuannya melalui publikasi pada forum konferensi internasional dan jurnal ilmiah bereputasi. Fokus dan minat penelitian beliau mencakup Computer Vision, Deep Learning, Vision Transformer, Facial Attribute Recognition, Natural Language Processing, Machine Learning, dan Algorithmic Fairness. Melalui kegiatan akademik dan eksplorasi riset yang konsisten, beliau terus memperdalam kompetensi keilmuan serta berkontribusi nyata dalam kemajuan teknologi komputasi cerdas. Rekam jejak penelitian serta profil akademik beliau dapat diakses melalui tautan \href{https://orcid.org/0009-0000-8021-1833}{ORCID: 0009-0000-8021-1833}.
\end{IEEEbiography}

\begin{IEEEbiography}[{\includegraphics[width=1in,height=1.25in,clip,keepaspectratio]{images/author_yuni.jpg}}]{Yuni Yamasari}
merupakan dosen dan peneliti di Departemen Informatika, Fakultas Informatika, Universitas Negeri Surabaya (UNESA) dengan landasan keilmuan yang kokoh di dunia akademik maupun praktis. Beliau menyelesaikan pendidikan Sarjana, Magister, hingga Doktor di Institut Teknologi Sepuluh Nopember (ITS), dengan kepakaran utama pada bidang Software Engineering dan Data Mining. Aktivitas tridarma beliau ditandai oleh keterlibatan aktif dalam pemberdayaan masyarakat, publikasi artikel ilmiah bereputasi, serta penulisan buku teks akademik. Beliau secara berkala mendiseminasikan temuan penelitiannya di berbagai forum ilmiah nasional maupun internasional dan telah memperoleh beberapa Hak Kekayaan Intelektual (HKI). Dedikasi akademik beliau diarahkan untuk mendorong perkembangan rekayasa perangkat lunak, analisis data komprehensif, dan bidang informatika terapan. Profil akademik dan riwayat publikasi beliau dapat ditelusuri melalui tautan \href{https://orcid.org/0000-0001-9719-3491}{ORCID: 0000-0001-9719-3491}.
\end{IEEEbiography}

\begin{IEEEbiography}[{\includegraphics[width=1in,height=1.25in,clip,keepaspectratio]{images/author_rafy.jpg}}]{Rafy Aulia Akbar}
menyelesaikan pendidikan Sarjana Teknik Informatika pada tahun 2019 dan meraih gelar Magister Informatika di Universitas Negeri Surabaya dengan konsentrasi pada bidang Ilmu Komputer. Semasa studi, beliau berhasil menorehkan prestasi dengan menerima penghargaan Best Solution Award pada kompetisi regional ASEAN Vocational and Engineering Camp 2017 bertema ``Toward School 4.0''. Saat ini, beliau aktif mendedikasikan diri sebagai peneliti di Departemen Informatika, Fakultas Informatika, Universitas Negeri Surabaya. Fokus dan minat penelitian beliau terpusat pada domain Computer Vision, Medical Image Analysis, Multimodal Deep Learning, Natural Language Processing, serta arsitektur Vision Transformer, khususnya dalam pengembangan sistem pembangkitan laporan medis otomatis. Rekam jejak publikasi dan profil ilmiah beliau dapat diakses melalui tautan \href{https://orcid.org/0009-0003-6991-0694}{ORCID: 0009-0003-6991-0694}.
\end{IEEEbiography}
